\section{通用逼近定理说清了什么又没说什么}
\label{sec:14-Deep-Learning/01-Representation-and-Approximation/02}


你刚搭好一个带 ReLU 隐藏层的网络，在 MNIST 上跑出了 97\% 的准确率。开心了五分钟之后，一个不安的念头冒出来：这个网络到底能表示什么样的函数？如果目标函数恰好落在它的表示范围之外，再好的优化器、再多的数据也救不回来。

第一篇已经回答了一半：没有非线性，再深的网络也只是线性映射，连 XOR 都分不开。插进激活函数，网络才能弯曲决策边界。但``能弯曲''是个很模糊的说法——弯曲到什么程度？一层隐藏层够不够？需要多少神经元？

你可能会想：既然一层就能弯曲，那堆很多层是不是纯粹浪费？或者反过来想：是不是有些函数一层网络永远够不着，必须加深？

这些问题的标准答案通常用一句话打发：``通用逼近定理说，一层就够了。''然后你就被这句话安慰了，心安理得地继续调参。但这句话同时说清了什么、又刻意回避了什么，才是真正影响你判断的东西。本节把定理本身说准确，给出它为什么成立的直觉，然后逐条检查它确实没有告诉你的事情。定理值得认真对待，也值得被诚实限定——它只回答表示、优化、泛化三问中的第一问，而且只是第一问的一小部分。

\subsection{先感受一下一层网络到底能做什么}

取最简单的设置：输入是一维实数 \(x\in[0,1]\)，隐藏层用 ReLU，输出是它们的线性组合。一个 ReLU 单元的形状是一条折线：左侧为零，右侧以某个斜率上升，弯折点由 bias 控制。两个 ReLU 相减，可以拼出一个``凸起''——第一篇已经展示过，三个 hinge 拼出一个三角形局部特征。

这个构造暗示了一件事：如果你有足够多的 ReLU 单元，你就可以在 \([0,1]\) 上摆满凸起，每个凸起的高度独立可调。把凸起换成更一般的``台阶''（sigmoid 版本），或者把凸起不断缩窄、增高（ReLU 版本），你就能逼近任意形状的连续函数——至少在一维是这样。

但一到高维，事情立刻变了味道。输入是 \(x\in[0,1]^d\)，每个隐藏单元只依赖一个方向 \(w^{\trans}x\)。也就是说，一个单元只能沿某一方向``切一刀''，产生一个沿该方向变化的 ridge。要覆盖整个 \(d\) 维立方体，你需要沿各个方向切很多刀。切多少刀？直觉上，把 \([0,1]^d\) 切成边长 \(\varepsilon\) 的小立方体，格子总数是 \((1/\varepsilon)^d\)。如果你只能靠逐格摆放小台块来逼近一个没有特殊结构的连续函数，那宽度至少要和格子数同一个量级。

这不是一个精确下界，但它已经把紧张感传递出来了：维数稍微上去一点，``一层就够''这句话在计算上就可能毫无意义。

\subsection{定理：一层非线性网络在紧集上稠密}

先把直觉收拢成精确陈述。固定激活函数 \(\phi\)，考虑所有单隐层前馈网络给出的函数集合

\[
\mathcal H_\phi
=\set{\,x\mapsto\sum_{j=1}^{m}\alpha_j\,
\phi\bigl(w_j^{\trans} x+b_j\bigr)
:\ m\in\N,\ \alpha_j\in\R,\ w_j\in\R^d,\ b_j\in\R\,}.
\]

通用逼近定理断言：只要 \(\phi\) 连续且非常数多项式——sigmoid、tanh、ReLU、GELU 全部满足——\(\mathcal H_\phi\) 在紧集上的连续函数空间里按 sup 范数稠密。写成可以直接使用的形式：

对任意连续函数 \(f:[0,1]^d\to\R\) 和任意精度 \(\varepsilon>0\)，存在某个宽度 \(m\) 和一组参数，使得

\[
\sup_{x\in[0,1]^d}
\abs{f(x)-\sum_{j=1}^{m}\alpha_j\,
\phi\bigl(w_j^{\trans} x+b_j\bigr)}
<\varepsilon.
\]

每个前提都值得停顿一下。

激活函数非多项式是必要条件：多项式的线性组合仍是多项式，最高次数被 \(\phi\) 的次数封死，无论如何叠加都跳不出多项式空间，而连续函数空间比多项式大得多。sigmoid 和 tanh 有饱和区，ReLU 是分段线性，GELU 光滑且非多项式——它们各不相同，但都满足``不是多项式''这个共同底线。

定义域取紧集同样不能省略。连续函数在无界域 \(\R^d\) 上可以发散到无穷，sup 范数逼近根本无从定义。锁定在 \([0,1]^d\) 或任何有界闭集上，连续函数自动一致连续且有界，逼近论的工具箱才能打开。

误差用 sup 范数度量，是对逼近质量的最强要求：它要求网络输出和目标函数在定义域内每一点的差距都不超过 \(\varepsilon\)。这个要求本身就排除了带间断的目标——在跳跃点附近，无论网络输出取什么值，它至少和跳跃两侧之一的距离超过跳跃幅度的一半，sup 误差压不下去。想处理间断函数，只能降级为积分范数（比如 \(L^2\)），相应的``逼近''含义也会改变。

一个熟悉的间断例子是购物满减规则：``消费超过 \(c\) 元免运费''。写成 \(a(x)=\ind\{x>c\}\)，金额只多一分钱，结果从不免跳成免。连续网络可以把过渡压缩到很窄的金额区间里——比如让 sigmoid 的斜率极大，使得 \(x<c-\delta\) 时输出几乎为零、\(x>c+\delta\) 时输出几乎为一——却永远不能在包含 \(c\) 的一个开区间上把 sup 误差压到任意小。这不是网络不够宽的问题，而是连续函数和间断函数之间的鸿沟，宽度填不平。

设备保护逻辑沿用的是同一条边界：实际系统通常让网络估计连续量（温度趋势、故障风险、安全裕度），再由单独的决策规则处理阈值、迟滞和误报代价。通用逼近定理担保的是前一个连续量的表示能力，不是硬停机策略的安全证明。

\subsection{直觉证明：先拼台阶，再拼函数}

严格证明需要 Hahn--Banach 和 Riesz 表示定理这类泛函分析工具，但核心画面在一维就能完整呈现。

取 sigmoid \(\sigma(z)=1/(1+e^{-z})\)。把它左右平移、拉缩放陡：

\[
g(x)=\sigma\bigl(a(x-x_0)\bigr)-\sigma\bigl(a(x-x_1)\bigr),
\qquad x_0<x_1.
\]

当 \(a\) 很大时，第一项在 \(x_0\) 左侧几乎为 0、右侧几乎为 1，第二项在 \(x_1\) 处把这个``开灯''动作关掉。两项相减，\(g\) 近似一个在 \([x_0,x_1]\) 上取 1、其余地方取 0 的方波——一个矩形台阶。

有了台阶，就可以搭楼梯。在 \([0,1]\) 上等距摆放 \(N\) 个互不重叠的窄台阶，每个台阶的高度独立可调，你就得到了任意阶梯函数。隐层单元的宽度 \(m\) 越大，可摆放的台阶越多、越窄。

剩下的一步是经典分析：紧集上的连续函数一致连续。把 \([0,1]\) 切成 \(N\) 等份，每份上取函数在左端点的值作为常数近似，得到的阶梯函数与原函数的 sup 误差不超过一致连续性给出的界，且随 \(N\) 增大而趋于零。两步拼接：连续函数先被阶梯逼近，阶梯再被 sigmoid 叠加逼近，总误差不超过两阶段误差之和。令 \(N\) 和 \(a\) 同时趋于无穷（宽度相应地增长），总误差趋于零。

ReLU 版本不需要新画面。ReLU 本身是 hinge：\(\max(0,z)\)。两个平移后的 hinge 相减（再加一个调整）产生一个三角形凸起。足够多的三角形凸起可以拼出任意分段线性插值，而连续函数在紧集上同样能被分段线性函数一致逼近。区别只在于：sigmoid 用台阶逼近阶梯，ReLU 用三角形逼近折线——两组``原子函数''不同，但覆盖连续函数空间的能力相同。

高维情形多了技术细节，思路不变：沿各方向构造只依赖 \(w^{\trans}x\) 的 ridge 函数，再把这些 ridge 叠成多维``台块''。等价地说，非多项式激活经平移缩放后提供了足够丰富的``基底''，可以像 Fourier 级数用正弦余弦填满平方可积空间那样，填满连续函数空间。

这个直觉构造同时也暴露了一个事实：它给出的宽度 \(m\) 随精度 \(\varepsilon\) 和维数 \(d\) 的恶化是指数的。定理对这件事保持沉默。

\subsection{没说之一：需要多宽}

定理对宽度 \(m\) 没有任何保证。它只说存在某个有限 \(m\)，而这个 \(m\) 可以随精度和维数爆炸。上一节的拼台阶构造已经把原因摊在桌上：把 \([0,1]^d\) 切成边长 \(\delta\) 的网格，格子数是 \((1/\delta)^d\)；对任意连续函数，没有额外结构可用时，逐格逼近是唯一的通用手段。维数灾难没有因为这条定理而蒸发。

Barron 的经典结果提供了一个对照。若目标函数的 Fourier 变换满足某种一阶矩有界条件——粗略地说，函数的高频成分衰减得足够快——那么单隐层网络的平方积分误差可以按 \(O(1/m)\) 的量级随宽度下降，速率中不出现维数 \(d\)。这说明：对那些``结构良好''的函数，确实存在可控的宽度保证。定理的普适性恰恰来自它对目标函数不假设任何结构——对全体连续函数成立，就意味着对最坏的那个连续函数只能给出存在性，不能给出有效宽度。

``可逼近''和``高效逼近''是两件不同的事，这也是整个表示论的一条分界线。

\subsection{没说之二：存在不等于找得到}

定理断言的是参数的存在性，证明手段是泛函分析中的稠密性论证和反证法，全程没有涉及任何算法。梯度下降面对的是一个非凸目标：上一单元已经看到，训练是数学目标、随机优化和数值实现共同形成的动态系统。步长、初始化、mini-batch 噪声、饱和区、梯度消失或爆炸——其中没有一环由逼近定理担保。

第一篇末尾还指出了参数空间的对称性：交换隐藏层两个神经元的编号，或者同时翻转一组权重的符号（对奇函数激活），网络函数不变。同一个函数因此被大量相距很远的参数表示。存在性证明既不利用这些简并，也不负责告诉你梯度下降会掉进哪一个参数盆地。

换句话说，定理给出的是``搜索空间里确实埋着解''，训练要回答的却是``从随机初始点出发、用随机梯度信号、在有限步数内能不能走到一个好解''。两个命题之间的跨度覆盖了几乎整个优化理论。

\subsection{没说之三：拟合不等于泛化}

即使宽度负担得起、优化也奇迹般地成功了，训练能检验的仍然只有有限个样本点上的函数值。一个足够宽的单隐层网络可以把训练数据拟合到任意精度——包括把纯随机标签也完美拟合——同时在训练点之间取几乎任意的值。插值本身不提供点之间的信息。

这个现象在本书前面已经多次出现。\hyperref[ch:063]{``插值与逼近：通过数据点，不等于理解点之间''}讨论过完全相同的困境：经过 \(n\) 个数据点的函数有无穷多个，训练损失只挑出了其中一个，对点之间的行为没有任何发言权。\hyperref[ch:118]{``统计学习理论：训练误差对真风险有多乐观''}把它量化：假设类的容量需要用样本量付费，训练误差对真风险存在系统性的乐观偏差。\hyperref[sec:13-Machine-Learning/10-Learning-Theory/04]{``泛化界的解释力与边界''}进一步指出，过参数化网络甚至能以零训练误差记住纯随机标签——能拟合训练集这一点，对泛化几乎什么也没有说明。

通用逼近定理告诉你网络``有足够容量''。但容量是把双刃剑：容量大到能容纳真实规律，也能大到把噪声原样记住。定理只说容量存在，不管你怎么使用它。

\subsection{深度：用复合换宽度}

回到开头那个疑问：既然一层就稠密，为什么还需要深度？

第一篇已经埋下了线索。浅层网络把 hinge 平铺在输入空间里，每个隐藏单元独立地在输入空间上切一刀。深层网络则不同：第二层单元作用在第一层输出的组合上，相当于在已经加工过的特征上再加工。把那个三角形映射复合 \(k\) 次，输出的折点数随 \(k\) 指数增长；浅层网络要在输入空间直接铺出同样多的折点，宽度必须指数增加，而深度方向只需线性加层。

这件事有正式的宽度--深度分离定理支撑。对天然具有复合结构的函数——先由局部边缘拼成部件轮廓，再由部件拼成物体——深层网络可以用多项式规模的参数和单元表示，浅层网络却需要指数宽度。注意：深度并没有带来浅层网络完全无法逼近的新函数类；它的价值在于高效逼近那些浅层``理论上也能逼近、实际上逼近不起''的函数。

但这个高效不是免费的。层数加深让优化景观更曲折——梯度要穿过多层复合，每一步都可能遭遇饱和、消失或爆炸。表示效率的提高以优化难度为代价。定理对这一切继续保持沉默，因为它的论域里根本没有算法。

\subsection{落点：表示、优化、泛化是三件不同的事}

把学习问题拆成三个独立的问题，比任何一条定理都更有助于判断什么时候该担心什么：

\begin{itemize}
\tightlist
\item
  表示：假设类中是否存在一个足够接近目标函数的候选；
\item
  优化：训练算法能否在有限资源内找到这样的参数；
\item
  泛化：找到的函数在训练分布之外是否仍然成立。
\end{itemize}

通用逼近定理只回答第一问，且只到``存在某个宽度''为止。宽度效率和深度分工属于第一问的精细部分，定理同样不提供答案。第二问交给本单元后续三篇：损失函数定义``找到什么算好''，反向传播提供梯度，训练动力学研究梯度能否真的把参数带过去。第三问整体移交\hyperref[ch:118]{``统计学习理论：训练误差对真风险有多乐观''}。

把三问混为一谈，是``通用逼近''被误用的主要方式。它从来不是``神经网络必然学得好''的免检证书，只是``表示空间没有硬天花板，值得动手训练''的入场券。入场之后，困难才真正开始。
